\chapter{MEMÓRIA DAS MEDIÇÕES DAS CICATRIZES E SELEÇÃO DOS PERFIS REPRESENTATIVOS}\label{ap:medicoes-cicatrizes}
\setcounter{figure}{0}\renewcommand{\thefigure}{C.\arabic{figure}}
\setcounter{table}{0}\renewcommand{\thetable}{C.\arabic{table}}

Este apêndice apresenta os valores utilizados na caracterização morfométrica das cicatrizes associadas ao MTD 1, ao MTD 2, ao evento C do MTC 2 e às duas cicatrizes do evento D do MTC 2. Para cada cicatriz, foram comparados perfis posicionados em diferentes setores da escarpa.

A altura da cicatriz (Hs) foi calculada a partir da diferença de tempo duplo de trânsito entre o topo e a base da escarpa, representados pelos pontos Aₛ e Bₛ, utilizando velocidade de propagação de 1.500 m/s. A inclinação máxima da cicatriz (Sₛ) foi calculada no segmento de maior declividade contínua de cada perfil, delimitado pelos pontos P e Q.

A seleção do perfil representativo foi baseada no maior valor de inclinação Sₛ. Após a seleção, a altura Hs adotada para a cicatriz correspondeu à altura medida nesse mesmo perfil. Portanto, os valores finais de Hs e Sₛ foram sempre obtidos em um único perfil, evitando a combinação de parâmetros provenientes de diferentes posições da escarpa.

\section{Cicatriz do MTD 1}

A cicatriz do MTD 1 foi avaliada por meio dos perfis esquerdo, central e direito. O perfil esquerdo apresentou a maior inclinação, de 12,36°, e foi adotado como representativo. A altura correspondente a esse perfil foi de 30,26 m (Tabela~\ref{tab:resultado-C-1}).

\begin{table}[H]
  \centering
  \caption{Medidas obtidas nos perfis avaliados na cicatriz do MTD 1.}
  \label{tab:resultado-C-1}
  \scriptsize
  \setlength{\tabcolsep}{2.5pt}
  \begin{tabularx}{\textwidth}{>{\raggedright\arraybackslash}X >{\centering\arraybackslash}X >{\centering\arraybackslash}X >{\centering\arraybackslash}X >{\centering\arraybackslash}X >{\centering\arraybackslash}X >{\centering\arraybackslash}X >{\centering\arraybackslash}X}
    \toprule
    \textbf{Perfil} & \textbf{ΔTWT Aₛ--Bₛ (ms)} & \textbf{Distância Aₛ--Bₛ (m)} & \textbf{Hs (m)} & \textbf{ΔTWT P--Q (ms)} & \textbf{Distância P--Q (m)} & \textbf{Gradiente (ms/km)} & \textbf{Sₛ (°)} \\
    \midrule
    Esquerdo — selecionado & 40,350 & 220,913 & 30,26 & 16,141 & 55,228 & 292,25 & 12,36 \\
    \midrule
    Central & 26,597 & 219,341 & 19,95 & 9,357 & 54,835 & 170,64 & 7,29 \\
    \midrule
    Direito & 29,128 & 211,534 & 21,85 & 9,590 & 52,883 & 181,42 & 7,75 \\
    \bottomrule
  \end{tabularx}
\end{table}

\nota{o perfil esquerdo foi selecionado por apresentar o maior valor de Sₛ. Os valores adotados para a cicatriz foram Hs = 30,26 m e Sₛ = 12,36°.}

\section{Cicatriz do MTD 2}

A cicatriz do MTD 2 também foi avaliada por meio de três perfis. O perfil direito apresentou a maior inclinação, de 9,79°, e foi selecionado como representativo. A altura correspondente a esse perfil foi de 16,23 m (Tabela~\ref{tab:resultado-C-2}).

\begin{table}[H]
  \centering
  \caption{Medidas obtidas nos perfis avaliados na cicatriz do MTD 2.}
  \label{tab:resultado-C-2}
  \scriptsize
  \setlength{\tabcolsep}{2.5pt}
  \begin{tabularx}{\textwidth}{>{\raggedright\arraybackslash}X >{\centering\arraybackslash}X >{\centering\arraybackslash}X >{\centering\arraybackslash}X >{\centering\arraybackslash}X >{\centering\arraybackslash}X >{\centering\arraybackslash}X >{\centering\arraybackslash}X}
    \toprule
    \textbf{Perfil} & \textbf{ΔTWT Aₛ--Bₛ (ms)} & \textbf{Distância Aₛ--Bₛ (m)} & \textbf{Hs (m)} & \textbf{ΔTWT P--Q (ms)} & \textbf{Distância P--Q (m)} & \textbf{Gradiente (ms/km)} & \textbf{Sₛ (°)} \\
    \midrule
    Esquerdo & 41,501 & 301,391 & 31,13 & 9,722 & 60,278 & 161,29 & 6,90 \\
    \midrule
    Central & 28,746 & 174,170 & 21,56 & 12,551 & 58,057 & 216,18 & 9,21 \\
    \midrule
    Direito — selecionado & 21,644 & 115,901 & 16,23 & 13,326 & 57,951 & 229,96 & 9,79 \\
    \bottomrule
  \end{tabularx}
\end{table}

\nota{o perfil direito foi selecionado por apresentar o maior valor de Sₛ. Os valores adotados para a cicatriz foram Hs = 16,23 m e Sₛ = 9,79°. Os valores desta tabela substituem as medições preliminares anteriormente obtidas para o perfil direito.}

\section{Cicatriz do evento C do MTC 2}

A cicatriz associada ao evento C do MTC 2 foi anteriormente identificada nos arquivos de trabalho como “cicatriz 3 de Búzios”. Entre os perfis cujos valores foram consolidados, o perfil esquerdo apresentou a maior inclinação, de 4,80°, e foi adotado como representativo. A altura correspondente foi de 32,85 m (Tabela~\ref{tab:resultado-C-3}).

\begin{table}[H]
  \centering
  \caption{Medidas obtidas nos perfis consolidados para a cicatriz do evento C do MTC 2.}
  \label{tab:resultado-C-3}
  \scriptsize
  \setlength{\tabcolsep}{2.5pt}
  \begin{tabularx}{\textwidth}{>{\raggedright\arraybackslash}X >{\centering\arraybackslash}X >{\centering\arraybackslash}X >{\centering\arraybackslash}X >{\centering\arraybackslash}X >{\centering\arraybackslash}X >{\centering\arraybackslash}X >{\centering\arraybackslash}X >{\centering\arraybackslash}X >{\centering\arraybackslash}X}
    \toprule
    \textbf{Perfil} & \textbf{Posição de Aₛ (m)} & \textbf{Posição de Bₛ (m)} & \textbf{ΔTWT Aₛ--Bₛ (ms)} & \textbf{Distância Aₛ--Bₛ (m)} & \textbf{Hs (m)} & \textbf{ΔTWT P--Q (ms)} & \textbf{Distância P--Q (m)} & \textbf{Gradiente (ms/km)} & \textbf{Sₛ (°)} \\
    \midrule
    Esquerdo — selecionado & 303,630 & 809,679 & 43,803 & 506,050 & 32,85 & 22,667 & 202,420 & 111,98 & 4,80 \\
    \midrule
    Central & 683,737 & 1.056,685 & 17,882 & 372,948 & 13,41 & 7,986 & 124,316 & 64,24 & 2,76 \\
    \bottomrule
  \end{tabularx}
\end{table}

\nota{as posições de Aₛ e Bₛ correspondem às distâncias acumuladas ao longo de cada perfil, medidas a partir do início do traçado. O perfil esquerdo foi selecionado por apresentar o maior valor de Sₛ. Os valores adotados para a cicatriz foram Hs = 32,85 m e Sₛ = 4,80°.}

\section{Cicatrizes do evento D do MTC 2}

O evento D do MTC 2 apresenta duas cicatrizes principais. As cicatrizes foram avaliadas separadamente, mantendo-se para cada uma o perfil com maior inclinação.

\subsection{Cicatriz 1}

Na cicatriz 1, o perfil direito apresentou a maior inclinação, de aproximadamente 7,30°, e foi adotado como representativo. A altura correspondente foi de aproximadamente 69,80 m (Tabela~\ref{tab:resultado-C-4}).

\begin{table}[H]
  \centering
  \caption{Medidas obtidas nos perfis avaliados na cicatriz 1 do evento D do MTC 2.}
  \label{tab:resultado-C-4}
  \scriptsize
  \setlength{\tabcolsep}{2.5pt}
  \begin{tabularx}{\textwidth}{>{\raggedright\arraybackslash}X >{\centering\arraybackslash}X >{\centering\arraybackslash}X >{\centering\arraybackslash}X >{\centering\arraybackslash}X >{\centering\arraybackslash}X >{\centering\arraybackslash}X >{\centering\arraybackslash}X}
    \toprule
    \textbf{Perfil} & \textbf{ΔTWT Aₛ--Bₛ (ms)} & \textbf{Distância Aₛ--Bₛ (m)} & \textbf{Hs (m)} & \textbf{ΔTWT P--Q (ms)} & \textbf{Distância P--Q (m)} & \textbf{Gradiente (ms/km)} & \textbf{Sₛ (°)} \\
    \midrule
    Esquerdo & 97,644 & 1.499,371 & 73,23 & 5,401 & 55,532 & 97,27 & 4,17 \\
    \midrule
    Central & 38,571 & 492,418 & 28,93 & 5,947 & 54,713 & 108,70 & 4,66 \\
    \midrule
    Direito — selecionado & ≈ 93,073 & ≈ 1.051,879 & ≈ 69,80 & 10,792 & 63,210 & 170,74 & 7,30 \\
    \bottomrule
  \end{tabularx}
\end{table}

\nota{o perfil direito foi selecionado por apresentar o maior valor de Sₛ. Os valores adotados para a cicatriz 1 foram Hs ≈ 69,80 m e Sₛ = 7,30°. O sinal de aproximação foi mantido porque os limites Aₛ--Bₛ dessa medição foram definidos aproximadamente no perfil.}

\subsection{Cicatriz 2}

Na cicatriz 2, o perfil esquerdo apresentou a maior inclinação, de 3,91°, e foi adotado como representativo. A altura correspondente foi de 40,60 m (Tabela~\ref{tab:resultado-C-5}).

\begin{table}[H]
  \centering
  \caption{Medidas obtidas nos perfis avaliados na cicatriz 2 do evento D do MTC 2.}
  \label{tab:resultado-C-5}
  \scriptsize
  \setlength{\tabcolsep}{2.5pt}
  \begin{tabularx}{\textwidth}{>{\raggedright\arraybackslash}X >{\centering\arraybackslash}X >{\centering\arraybackslash}X >{\centering\arraybackslash}X >{\centering\arraybackslash}X >{\centering\arraybackslash}X >{\centering\arraybackslash}X >{\centering\arraybackslash}X >{\centering\arraybackslash}X >{\centering\arraybackslash}X}
    \toprule
    \textbf{Perfil} & \textbf{Posição de Aₛ (m)} & \textbf{Posição de Bₛ (m)} & \textbf{ΔTWT Aₛ--Bₛ (ms)} & \textbf{Distância Aₛ--Bₛ (m)} & \textbf{Hs (m)} & \textbf{ΔTWT P--Q (ms)} & \textbf{Distância P--Q (m)} & \textbf{Gradiente (ms/km)} & \textbf{Sₛ (°)} \\
    \midrule
    Esquerdo — selecionado & 242,234 & 1.150,613 & 54,136 & 908,379 & 40,60 & 16,558 & 181,676 & 91,14 & 3,91 \\
    \midrule
    Central & 339,113 & 1.017,338 & 45,647 & 678,225 & 34,23 & 17,936 & 203,468 & 88,15 & 3,78 \\
    \midrule
    Direito & 512,228 & 922,011 & 26,089 & 409,783 & 19,57 & 12,879 & 153,668 & 83,81 & 3,60 \\
    \bottomrule
  \end{tabularx}
\end{table}

\nota{as posições de Aₛ e Bₛ correspondem às distâncias acumuladas ao longo de cada perfil. O perfil esquerdo foi selecionado por apresentar o maior valor de Sₛ. Os valores adotados para a cicatriz 2 foram Hs = 40,60 m e Sₛ = 3,91°.}

\section{Síntese dos perfis selecionados}

A síntese dos perfis selecionados é apresentada na Tabela~\ref{tab:resultado-C-6}.

\begin{table}[H]
  \centering
  \caption{Síntese dos perfis e parâmetros adotados na caracterização das cicatrizes.}
  \label{tab:resultado-C-6}
  \scriptsize
  \setlength{\tabcolsep}{2.5pt}
  \begin{tabularx}{\textwidth}{>{\raggedright\arraybackslash}X >{\centering\arraybackslash}X >{\centering\arraybackslash}X >{\centering\arraybackslash}X >{\centering\arraybackslash}X >{\centering\arraybackslash}X >{\centering\arraybackslash}X >{\centering\arraybackslash}X >{\centering\arraybackslash}X >{\centering\arraybackslash}X}
    \toprule
    \textbf{Depósito ou evento} & \textbf{Cicatriz} & \textbf{Perfil selecionado} & \textbf{ΔTWT Aₛ--Bₛ (ms)} & \textbf{Distância Aₛ--Bₛ (m)} & \textbf{Hs (m)} & \textbf{ΔTWT P--Q (ms)} & \textbf{Distância P--Q (m)} & \textbf{Gradiente (ms/km)} & \textbf{Sₛ (°)} \\
    \midrule
    MTD 1 & Única & Esquerdo & 40,350 & 220,913 & 30,26 & 16,141 & 55,228 & 292,25 & 12,36 \\
    \midrule
    MTD 2 & Única & Direito & 21,644 & 115,901 & 16,23 & 13,326 & 57,951 & 229,96 & 9,79 \\
    \midrule
    MTC 2-C & Única & Esquerdo & 43,803 & 506,050 & 32,85 & 22,667 & 202,420 & 111,98 & 4,80 \\
    \midrule
    MTC 2-D & Cicatriz 1 & Direito & ≈ 93,073 & ≈ 1.051,879 & ≈ 69,80 & 10,792 & 63,210 & 170,74 & 7,30 \\
    \midrule
    MTC 2-D & Cicatriz 2 & Esquerdo & 54,136 & 908,379 & 40,60 & 16,558 & 181,676 & 91,14 & 3,91 \\
    \bottomrule
  \end{tabularx}
\end{table}

\nota{Hs representa a altura da cicatriz, Aₛ--Bₛ delimita a escarpa, P--Q delimita o segmento de maior declividade contínua e Sₛ representa sua inclinação. Todos os desníveis foram convertidos utilizando velocidade de 1.500 m/s.}
